% Subsection: Systems of Subsets as Structured Objects  (notes stub; content authored later)
\subsection{Systems of Subsets as Structured Objects}

\begin{remark*}[Exposition]
The phrase ``system of sets'' marks a change of viewpoint.  A family of
subsets can be treated as a structured object when it is studied together with
the operations it is meant to tolerate.  Later sections will make this precise
through closure under finite, countable, and generated operations.
\end{remark*}

\begin{definition}[Set System]\label{def:set-system}
Let \(X\) be a set.  A \emph{set system on \(X\)} is a family
\(\mathcal{S}\subseteq\mathcal{P}(X)\) whose members are studied collectively
as the selected subsets of \(X\).
\end{definition}

\begin{remark*}[Standard quantified statement]
\[
  \mathcal{S}\text{ is a set system on }X
  \quad\Longleftrightarrow\quad
  \mathcal{S}\subseteq\mathcal{P}(X).
\]
\end{remark*}

\begin{remark*}[Interpretation]
The definition is intentionally broad.  A set system need not yet be closed
under any operation.  The point is to isolate the object that later receives
extra requirements: closed under complements, closed under finite unions,
closed under countable unions, generated by a smaller collection, or compared
with another system on the same base set.
\end{remark*}

\begin{remark*}[Exposition]
Many later spaces are built by choosing a set system with special closure
rules.  The later chapters will attach geometric, topological, or measurable
meaning to those choices.  This chapter first develops the common algebraic
language: what it means to select subsets, apply operations to selected
subsets, and close a family under the operations demanded by a theory.
\end{remark*}

\begin{remark*}[Exposition]
Properties of a single subset \(A\subseteq X\) are point-level properties:
which elements are in \(A\), whether \(A\) contains another subset, or whether
two subsets are equal.  Properties of a set system \(\mathcal{S}\) are
family-level properties: which subsets of \(X\) are admitted into
\(\mathcal{S}\), and whether operations on admitted subsets produce admitted
subsets again.
\end{remark*}

\begin{dependencies}
\begin{itemize}
  \item \hyperref[def:family-of-subsets]{Family of Subsets}
\end{itemize}
\end{dependencies}
